\subsection{The paper's actual finite-flat function norm and base change}
\label{subsec:paper-flat-norm-base-change}

\begin{theorem}[Actual paper fibre multiplicities and whole fibre points]
\label{spectrum:paper-flat-fibre-identification}
\lean{EnvelopeProperty.Berk.flatFibreMult_eq_spectrumFibreNormMult,
EnvelopeProperty.Berk.absoluteFibreEquiv,
EnvelopeProperty.Berk.absoluteFibreEquiv_apply_val}
\leanok
\uses{spectrum:nonetale-fibre-multiplicities,spectrum:trivial-field-bridge,
spectrum:relative-integral-completed-fibres}
Let $k$ be a field with trivial absolute value and let $A\to B$ be a finite
map of commutative $k$-algebras. The existing paper multiplicity
$\mathrm{flatFibreMult}_A(y)$ equals the generic spectrum multiplicity,
namely the actual $H(x)$-dimension of the localized completed tensor fibre
at the kernel corresponding to $y$, where $x=y|_A$. The old paper fibre
and the whole generic spectrum fibre over $x$ are genuinely equivalent;
on every point this equivalence simply forgets the trivial-base condition.
\end{theorem}
\begin{proof}
The old and generic completed residue fields and completed tensor fibre
maps agree on the actual point and coefficient maps. The existing old
multiplicity formula and the generic local-dimension formula therefore
give the same actual finite-dimensional localization, with no residue-only
or reduced-fibre substitution. Restrict the genuine paper/relative-spectrum
homeomorphism to the whole base fibre and compose with the actual relative
fibre forget equivalence. The latter equips each absolute point over $x$
with its forced coefficient values, since its restriction to $A$ is the
whole point $x$. Converting equality fibres to inverse images of singleton
sets gives the stated equivalence and its evaluation identity.
\end{proof}

\begin{theorem}[The actual paper norm is the generic weighted sum]
\label{spectrum:paper-flat-norm-identification}
\lean{EnvelopeProperty.Berk.extendFunction,
EnvelopeProperty.Berk.extendFunction_toSpectrum,
EnvelopeProperty.Berk.flatNormFun_eq_spectrumWeightedSum,
EnvelopeProperty.Berk.extendFunction_comap}
\leanok
\uses{spectrum:paper-flat-fibre-identification}
Let $u$ be an extended-real function on $\operatorname{Berk}_k(B)$.
Extend it to the absolute spectrum by zero outside the trivial-$k$ locus.
The existing $\mathrm{flatNormFun}(u)(x)$ is exactly the generic finite-fibre
weighted sum of this extension. This actual extension commutes with
pullback through every $k$-algebra homomorphism.
\end{theorem}
\begin{proof}
On an old point, the defining coefficient values satisfy the extension's
positive branch, which returns the original value of $u$. Reindex the
existing finite paper sum through the preceding actual fibre equivalence,
and replace each multiplicity by the proved generic dimension weight.
The natural-number-to-real-to-EReal coefficient in the old definition is
exactly the generic natural-number EReal coefficient.
For pullback, an algebra homomorphism preserves every coefficient element;
the trivial-base conditions are thus equivalent, and both branches of the
actual extension agree. No regularizability or boundedness of $u$ is used.
\end{proof}

\begin{theorem}[Actual finite-flat function norms under coefficient base change]
\label{spectrum:paper-flat-norm-base-change}
\lean{EnvelopeProperty.Berk.flatNormFun_baseChange,
EnvelopeProperty.Berk.flatNormFun_scalarExtension}
\leanok
\uses{spectrum:paper-flat-norm-identification,
spectrum:all-function-finite-base-change,spectrum:all-function-scalar-extension}
Assume only that $B$ is finite over $A$. For every commutative $k$-algebra
$C$ over $A$, every actual point $x'$ of $\operatorname{Berk}_k(C)$, and
every EReal function $u$ on $\operatorname{Berk}_k(B)$,
\[
 \mathrm{flatNorm}_{C\otimes_A B/C}(u\circ r_B)(x')
 =\mathrm{flatNorm}_{B/A}(u)(x'|_A).
\]
For an arbitrary commutative $k$-algebra $S$, the same identity holds for
the \emph{literal} scalar-extended map
$S\otimes_kA\to S\otimes_kB$, with the actual right-inclusion restriction
maps. Thus it applies in particular at supplied points with any prescribed
compatible absolute value on an extension field $S$; it does not require
the value on $S$ to be trivial or $S$ to be complete.
The statements hold for every EReal function with Mathlib's actual
extended-real additive conventions. Flatness, freeness, projectivity,
separability, reducedness, infinitude and upper-finiteness are unnecessary.
\end{theorem}
\begin{proof}
Rewrite both existing paper norms as the generic weighted sums and use
compatibility of the extension with the actual pullback. The proved
arbitrary-function finite-fibre comparison gives the first equality.
For the literal scalar extension, apply its proved genuine tensor
cancellation and the all-function scalar-extension comparison: the tensor
isomorphism preserves every original element, while the local-dimension
conservation theorem retains all nilpotent and inseparable thickness.
Both comparisons regroup the genuine local dimension weights, rather
than assuming the scalar-extension fibres are bijective. All EReal
functions are permitted because EReal is an additive commutative monoid
with the unconditional natural-scalar identity used by the generic
producer. The given $x'$ is retained throughout; no point-extension
existence theorem is asserted here.

This is the affine pointwise base-change identity for the actual norm
appearing in Fang--Han--Xia Lemma~5.13. The separate author's theorem
that finite-flat norms preserve regularizable functions supplies that
additional function-space statement; it is not an assumption in this
pointwise proof.
\end{proof}
